\section{Evidence on the variance model and the interval}\label{sec:evidence}

The evidence comes in four steps. A method study on a six-year Veneto panel shows
which wages determine precision and how the mover model behaves when the truth is
known. A bottleneck study shows what the $q=1$ interval adds. Registered simulations
through the Stata command test the interval a user actually receives. The comparison
with KSS follows in Section~\ref{sec:kss}.

\subsection{Which wages determine precision?}
The Veneto Worker History panel for 1996--2001 contains 5.16 million
person-years, 1.11 million workers and 86,647 firms after pruning. Its
666,090 stayers contribute 3.03 million person-years. Collapsing movers'
years to matches leaves 4.06 million deletion units. This six-year design
provides enough repeated observations to study where the variance inputs
come from, including for the variance of the sum.

\begin{table}[htbp]\centering
\caption{Estimates and sources of sampling variance on the six-year Veneto panel}
\label{tab:real}
\begin{tabular}{@{}lrrr@{}}
\toprule
 & Leave-out estimate & Standard error & Movers' share (\%)\\
\midrule
Worker variance & 0.0970 & 0.000218 & 97\\
Firm variance   & 0.0261 & 0.000196 & 100\\
Worker--firm covariance & 0.0134 & 0.000096 & 100\\
Variance of the sum & 0.1499 & 0.000095 & 77\\
\bottomrule
\end{tabular}
\par\smallskip\begin{minipage}{0.92\linewidth}\footnotesize
Shares refer to the first term of \eqref{eq:Vhat}, rounded to whole
percentages. These are method-study estimates under match deletion, from a prototype
that follows the algorithm of \fevc{}. Year effects and two age terms were removed and
treated as known. The study used exact stayer leverages, 1,600 projections for movers
and 16 draws for the variance calculation.
\end{minipage}
\end{table}

The standard errors are small, between 0.06\% and 0.75\% of the estimates
(Table~\ref{tab:real}). Stayers supply nearly a quarter of the first term for the
variance of the sum with variance estimates that need no model. The rest of that term,
and nearly all of it for the other three components, rests on the model for movers.

Among movers, a small fraction of rows can have a large influence on
precision. Just 1\% of movers' rows lie beyond the cap on the estimated
worker effect, yet they account for 26\% of the first term of the estimated
sampling variance for the sum. Figure~\ref{fig:profile} shows why allowing
variances to depend on wage levels helps. Average leave-out products are
high at both extremes of the worker-effect distribution, and they fall as
the firm effect rises. The fitted profile captures much of this shape.
It overpredicts the worker-effect decile averages in deciles 2--4 and
underpredicts them in deciles 7--9. Because an estimated effect and the
raw leave-out product use some of the same wages, these averages are an
imperfect benchmark. The actual mover variances remain unobserved, which is why the
next subsection turns to simulations in which they are known.

\begin{figure}[htbp]
\centering
\def\profalphastay{(-1.826,2.678) (-0.846,1.459) (-0.523,0.960) (-0.298,0.757) (-0.109,0.660) (0.070,0.628) (0.262,0.641) (0.499,0.662) (0.874,0.784) (1.939,1.184)}
\def\profalphafit{(-1.826,2.249) (-0.846,1.268) (-0.523,1.048) (-0.298,0.928) (-0.109,0.849) (0.070,0.789) (0.262,0.741) (0.499,0.705) (0.874,0.695) (1.939,1.057)}
\def\profalphamove{(-1.713,2.424) (-0.838,1.389) (-0.522,0.948) (-0.299,0.810) (-0.110,0.760) (0.068,0.730) (0.258,0.736) (0.495,0.800) (0.867,0.881) (1.926,1.081)}
\def\profpsistay{(-1.868,1.352) (-0.796,1.076) (-0.414,1.021) (-0.176,1.028) (0.009,0.952) (0.187,0.937) (0.360,0.976) (0.564,0.962) (0.894,0.848) (1.838,0.913)}
\def\profpsifit{(-1.868,1.323) (-0.796,1.055) (-0.414,1.027) (-0.176,1.015) (0.009,0.981) (0.187,0.975) (0.360,0.949) (0.564,0.942) (0.894,0.904) (1.838,0.888)}
\def\profpsimove{(-1.868,1.451) (-0.798,1.041) (-0.419,0.988) (-0.178,0.942) (0.014,0.882) (0.183,0.884) (0.359,0.866) (0.571,0.931) (0.899,0.955) (1.680,1.124)}
\def\figalphalo{(-1.683,17.859) (-0.838,5.629) (-0.524,3.553) (-0.300,2.909) (-0.111,2.704) (0.067,2.634) (0.257,2.890) (0.493,3.626) (0.863,5.183) (1.932,9.254)}
\def\figalphafit{(-1.683,15.960) (-0.838,6.384) (-0.524,4.410) (-0.300,3.438) (-0.111,2.889) (0.067,2.565) (0.257,2.419) (0.493,2.479) (0.863,3.139) (1.932,10.138)}
\def\figpsilo{(-1.957,9.523) (-0.798,5.408) (-0.421,4.579) (-0.178,4.275) (0.014,4.235) (0.182,3.768) (0.356,3.728) (0.570,3.977) (0.904,3.852) (1.680,4.436)}
\def\figpsifit{(-1.957,9.363) (-0.798,5.346) (-0.421,4.734) (-0.178,4.301) (0.014,4.021) (0.182,3.933) (0.356,3.806) (0.570,3.706) (0.904,3.635) (1.680,4.392)}

\begin{tikzpicture}
\begin{groupplot}[group style={group size=2 by 1, horizontal sep=1.3cm}, width=0.5\textwidth, height=0.36\textwidth,
  ymin=0, ymax=20, xmin=-2.4, xmax=2.4, ymajorgrids, grid style={gray!20},
  xlabel style={font=\footnotesize}, ylabel style={font=\footnotesize}, ticklabel style={font=\footnotesize},
  title style={font=\footnotesize}, legend style={font=\footnotesize, draw=none, fill=none, legend cell align=left}]
\nextgroupplot[title={By estimated person effect}, xlabel={standardised person effect (decile means)},
  ylabel={variance of a mover's row ($\times100$)}, legend style={at={(0.98,0.97)}, anchor=north east}]
\addlegendimage{only marks, mark=*, mark size=1.8pt, cmove}
\addlegendentry{leave-out estimates, average}
\addlegendimage{no marks, thick, cstay}
\addlegendentry{fitted by the model, average}
\addplot[cmove, only marks, mark=*, mark size=1.8pt, forget plot] coordinates {\figalphalo};
\addplot[cstay, thick, no marks, forget plot] coordinates {\figalphafit};
\nextgroupplot[title={By estimated firm effect}, xlabel={standardised firm effect (decile means)}]
\addplot[cmove, only marks, mark=*, mark size=1.8pt] coordinates {\figpsilo};
\addplot[cstay, thick, no marks] coordinates {\figpsifit};
\end{groupplot}
\end{tikzpicture}
\caption{Movers' rows by decile of the estimated effects: average leave-out estimate $\hat\sigma_i^2$ and average
fitted variance $\st_i^2$. The bin boundaries are deciles of stayers' estimated effects. The plotted observations are movers' rows in the six-year Veneto method study.}
\label{fig:profile}
\end{figure}


\subsection{What changes when the truth is known?}\label{sec:sim}
Holding the six-year network and effects fixed, a method study drew new wages
under ten variance patterns, including flat variances, skewed and fat-tailed errors,
and different wage profiles for movers and stayers. Each scenario had 200 draws, and
each draw ran the estimator as on the data: projected leverages for movers, the two
regressions with estimated effects in the profile, and simulated double sums.
Table~\ref{tab:sim} reports the ratio of the root-mean-square standard error to the
benchmark standard deviation, that is, the square root of the mean variance estimate
over the true variance, together with the coverage of the normal 95\% interval.

\begin{table}[htbp]\centering\small\setlength{\tabcolsep}{4.5pt}
\caption{Simulation on the six-year Veneto design, with the fitted model for movers}
\label{tab:sim}
\begin{tabular}{@{}rl cccc c cccc@{}}
\toprule
 & & \multicolumn{4}{c}{Standard error / benchmark} & & \multicolumn{4}{c}{Coverage of the 95\% interval}\\
\cmidrule(lr){3-6}\cmidrule(l){8-11}
 & Scenario & Worker & Firm & Cov. & Sum & & Worker & Firm & Cov. & Sum\\
\midrule
\tabinput{\srcbase unified/revision/tab_sim_model.tex}
\bottomrule
\end{tabular}
\par\smallskip\begin{minipage}{0.97\linewidth}\footnotesize
A ratio of 1 means that the mean variance estimate equals the true variance. 200 draws
per scenario; the Monte Carlo standard error of a coverage rate is about 0.015. In
scenarios 1 to 7 movers and stayers share one variance profile; in 8 to 10 movers have
their own. Scenarios 3 to 5 depart from the assumed shape of stayers' errors.
\end{minipage}
\end{table}

\emph{The repeat-measurement part is exact.} With the movers' true variances supplied,
the calculation isolates the contribution from repeated wages. In the eight scenarios
without skewness the ratio was 1.00--1.01 for every component, fat tails included.
With skewed stayer errors the ratio for the variance of the sum was 0.97 or 1.03,
depending on the sign of skewness relative to wages. These are the signs predicted by
the calculation in Appendix~\ref{app:proofs}.

\emph{With the fitted model the standard error stays close.} Across the forty
component--scenario cases the ratio lies between 0.96 and 1.10. It is never more than
4\% too small, and coverage is between 93.5\% and 98.0\%. This includes scenario 10,
in which movers follow the decile profile of their own leave-out products in the data
and which is not within the model: its ratios are 1.06 to 1.10.

\emph{Why the ratios lean above one.} Looking only at the ratios would conceal two
opposing contributions. With true effects and true variances as inputs to the fit,
misspecification alone would understate the standard error for the variance of the sum
by 4--8\% in eight scenarios. One source is the cap: where true variances keep rising
beyond three standard deviations, the 1\% of rows beyond the cap can carry a third of
the first term, and raising the cap from three to six standard deviations would move
the predicted ratio in that scenario from 0.95 to 0.99. Using estimated effects and
fitting on noisy wages, together with the other steps of the calculation, moved the
result in the opposite direction, often by more. The comparison cannot attribute the
difference to any one step, and no result predicts its sign in other designs.

\begin{table}[htbp]\centering\small\setlength{\tabcolsep}{4pt}
\caption{Validation with tolerances fixed in advance: ratio and coverage by design}
\label{tab:val}
\begin{tabular}{@{}lc cccc cc@{}}
\toprule
 & & \multicolumn{4}{c}{Standard error / benchmark (range)} & & \\
\cmidrule(lr){3-6}
Design & Scenarios & Worker & Firm & Cov. & Sum & Coverage & Met\\
\midrule
\tabinput{\srcbase note/gen/tab_val.tex}
\bottomrule
\end{tabular}
\par\smallskip\begin{minipage}{0.97\linewidth}\footnotesize
200 fresh draws per scenario. Ranges are over scenarios; coverage is the smallest and
largest value over scenarios and components. \emph{Met}: whether the design met every
registered tolerance. Registered tolerances: for every scenario and component a
ratio between 0.93 and 1.15 and coverage of at least 0.90; for every design an average
coverage of at least 0.94, no failed draw, no nonpositive variance estimate, and
floors binding for at most 2\% of movers' rows in any draw.
\end{minipage}
\end{table}

\emph{The results hold on fresh draws and on two other designs.} The ten scenarios
also informed the choice of model. The estimator was then frozen and run once more
with tolerances written down in advance, on fresh draws of the six-year design, on a
random half of its workers, and on a three-year window (Table~\ref{tab:val}). The
six-year and half-worker runs met every tolerance, and the six-year ratios reproduced
those of Table~\ref{tab:sim} to within 0.007. In the three-year window the
firm-variance standard error was 1.152 and 1.162 times its benchmark in two scenarios,
above the registered limit of 1.15. Across the 88 cases of the three designs the ratio
was below 0.95 in one and above 1.10 in eleven. Shorter panels give less information
about each worker's effect, and there the standard error of firm variance errs on the
conservative side by up to a sixth.

\subsection{What the \texorpdfstring{$q=1$}{q=1} interval adds at a bottleneck}\label{sec:q1sim}
A separate study built synthetic panels with a bottleneck of chosen width: 100 firms
in two groups of 50, 12,000 workers observed for two to six years, and $b$ bridging
workers who move once between the groups. The link has $b=100$, 30 or 10 workers, or
there is no bottleneck. Error variances are flat, follow a wage profile, are steeper
for movers, or are three times as high for the bridging workers alone, a feature the
mover model cannot see. The designs, 1,000 draws per design, and the tolerances were
fixed beforehand. The study tests the $q=1$ construction of
Section~\ref{sec:q1int} applied to every draw, not the automatic choice.

\begin{table}[htbp]\centering\footnotesize\setlength{\tabcolsep}{4.5pt}
\caption{Coverage of 95\% intervals with a bottleneck of $b$ bridging workers}
\label{tab:q1}
\begin{tabular}{@{}l c c c c c c c c@{}}
\toprule
 & & & \multicolumn{4}{c}{Coverage (smallest to largest)} & \multicolumn{2}{c}{Length of $q=1$} \\ \cmidrule(lr){4-7}\cmidrule(l){8-9}
Bottleneck & $\max_iw_{i1}^2$ & $\kappa$ & Normal & $q=1$ & Leave-out & True $\Sigma$ & /normal & /true $\Sigma$ \\
\midrule
None & .157--.167 & 0.05--0.28 & .939--.964 & .941--.970 & .941--.967 & .945--.963 & 1.04 & 1.01 \\
$b=100$ & .010 & 0.30--1.42 & .708--.956 & .951--.986 & .944--.984 & .953--.982 & 1.28 & 1.02 \\
$b=30$ & .031 & 0.95--5.09 & .528--.935 & .973--.993 & .935--.986 & .966--.990 & 1.60 & 1.08 \\
$b=10$ & .090 & 2.74--11 & .487--.937 & .970--1.00 & .795--.911 & .971--.996 & 1.90 & 1.19 \\
\addlinespace
$b=30,10$, skewed errors & .031--.090 & 1.11--9.90 & .743--.923 & .972--.995 & .812--.967 & .977--.993 & 1.60 & 1.08 \\
$b=30,10$, noisy bridge & .031--.090 & 2.77--23 & .629--.727 & .950--.991 & .921--.991 & .982--.996 & 2.25 & 1.00 \\
$b=3$ (stress test) & .256 & 11--35 & .713--.805 & .940--.980 & .710--.833 & .939--.989 & 2.16 & 1.11 \\
\bottomrule
\end{tabular}

\par\smallskip\begin{minipage}{0.97\linewidth}\footnotesize
1,000 draws per design; ranges over designs and over worker variance, firm variance
and the covariance. \emph{Normal}: $\hat\theta\pm1.96\,\widehat{\V}[\hat\theta]^{1/2}$.
\emph{$q=1$}: the interval of Section~\ref{sec:q1int} with \eqref{eq:q1est}.
\emph{Leave-out}: the same with KSS's $V_b^{\rm lo}$ alone. \emph{True $\Sigma$}: the
interval with the true covariance matrix, which no estimate can use. A draw without an
interval counts as not covered. \emph{Length}: mean length of the $q=1$ interval over
that of the normal interval, and over that of the interval with the true matrix.
$\max_iw_{i1}^2$ is the largest squared weight of one row in the leading direction and
$\kappa$ the curvature. Noisy bridge: the bridging workers' errors have three times
the variance.
\end{minipage}
\end{table}

\emph{The $q=1$ interval covers throughout.} All 27,000 draws produced a $q=1$
interval. Across the 75 registered component--design cases it covered between 94.1\%
and 100\%, on average 97.8\%, including with skewed errors and with noisy bridging
workers (Table~\ref{tab:q1}). It is conservative, and most of that is the construction
of KSS and not the estimates: with the true covariance matrix coverage is 94.5\% to
99.6\%.

\emph{The normal interval fails as the bottleneck narrows.} Its coverage falls to
48.7\% with ten bridging workers. The estimate is then far from normal, and the
variance estimate is nonpositive in up to 47.8\% of the draws. Where bridging workers
are noisier than the model can know, the normal standard error is 0.59 to 0.67 times
the true one. For the variance of the sum, which has no dominant direction, the normal
standard error is 1.00 to 1.06 times the true one in every design.

\emph{The larger-of-two rule earns its place.} With ten bridging workers KSS's
leave-out variance of the score alone gives coverage of 79.5--91.1\%, because the
products of a few rows are too noisy. With noisy bridging workers the fitted variance
alone is about a third of the true value and its interval covers 80.3--92.5\%. The
larger of the two covers in both situations.

One registered tolerance was missed. Without a bottleneck the $q=1$ interval was to be
at most 10\% longer than the normal one, and for the covariance in one design it was
12.7\% longer. The interval is not wrong there: it covered 96.9\% against 94.9\% for
the normal one. The result stands as a failure of the tolerance. A stress test with
only three bridging workers lies outside the coverage result
(Appendix~\ref{app:evidence}).

\subsection{Coverage of the interval selected by the command}\label{sec:impl-mc}
The coverage question for a user is whether the interval actually reported
contains the true component. A draw that yields no interval counts as a miss.
The package experiments applied the Stata command to two-block networks of
about 48,000 person-years, ranging from well connected to a single pronounced
bottleneck. They included both deletion modes, normal and skewed errors, and
variance patterns both within and outside the fitted model. Each confirmation
used 20 designs, two deletion modes and 1,000 replications per combination, which
gives 160 cases of a design, a deletion mode and a component.

\begin{table}[htbp]\centering
\caption{Registered confirmations of nominal 95\% intervals through the command}
\label{tab:impl-mc}
\begin{tabular}{@{}lrrrl@{}}
\toprule
Confirmation & Fits & Lowest coverage & Mean coverage & Outcome\\
\midrule
Without controls, first & 40,000 & 0.937 & 0.960 & Failed inventory\\
Without controls, second & 40,000 & 0.930 & 0.959 & Passed\\
With controls & 40,000 & 0.923 & 0.960 & Passed\\
\bottomrule
\end{tabular}
\par\smallskip\begin{minipage}{0.92\linewidth}\footnotesize
Coverage is over replications within a case; lowest and mean are over the 160 cases.
Registered tolerances: coverage of at least 0.92 in every case and 0.94 on average,
an interval for all four components in at least 99\% of the replications of a design
and deletion mode, and, without a bottleneck, a reported interval of median length at
most 1.10 times the normal one with \code{e(V)} posted in at least half of the
samples. In the first run one fit of 40,000 returned point estimates only. The
validation program of the time had no rule for that documented fallback and rejected
its task, so the registered verdict is a failure and stands; counting the fit as a
miss, every tolerance was met, and the coverage figures shown are descriptive. The
no-control confirmations used build \code{fb49e16d} and the controls confirmation
\code{f3938986}. These are pre-release builds; Appendix~\ref{app:evidence} describes
how far the evidence carries to the released package.
\end{minipage}
\end{table}

\emph{The reported interval covers.} In the two passed confirmations none of the
80,000 fits failed or fell back to point estimates. Over their 320 cases the reported
95\% interval covered 96.0\% on average, between 92.3\% and 99.6\% case by case. It
covered at least 95\% in 239 cases and less than 94\% in 16.

\emph{How to read the lowest case.} With 1,000 replications a coverage rate has a
sampling standard error of about 0.007. If every case had exactly 95\% coverage, the
lowest of 160 independent cases would typically be 0.931, and 0.935 if only 40 were
independent, as when the four components of a fit move together. The lowest case
without controls, 0.930, is therefore what nominal coverage looks like. The lowest case
with controls, 0.923, is a real but small shortfall, discussed below.

\emph{The switch is what delivers this.} On the second no-control confirmation, using
normal intervals throughout would have given a lowest coverage of 78.9\% on the same
replications, compared with 93.0\% for the selected interval. The switch costs little
where it is not needed: without a bottleneck the selected interval was at most 1.059
times as long as the normal one, and the command returned a joint covariance in 89\%
of those samples. Selecting the 10\% threshold involved a trade-off. In the development
runs, lowering it to 5\% raised the lowest coverage from 93.7\% to 94.0\%, but more
than doubled the frequency of selecting $q=1$ without a bottleneck, from 4.3\% to
9.8\%.

\emph{Controls.} The controls confirmation used new seeds and four kinds of controls.
Its lowest coverage occurred where the two blocks differed and a control loaded on the
worker effect. The worker- and firm-variance standard errors were 8\% and 7\% below
the spread of the estimates there. Separate development stress tests show why the
distinction between joint and fixed-offset estimation matters. In one case without a
bottleneck, under match deletion and with a control index as large as the effects, the
joint estimate's covariance-component standard error was 12\% too small. With a fixed
offset on the same design, the component standard errors ranged from 0.99 to 1.07 of
the spread. By registered design the controls confirmation excluded this large-index
case, variances that move with a control inside a cell, and a control loaded on the
firm effect when the blocks differ.

\section{Comparison with KSS}\label{sec:kss}\label{kss:sec:empirical}
\subsection{The same point estimator, three sources of variance information}
The split-sample procedure in KSS (2020) allows arbitrary heteroskedasticity.
It predicts the mean of an observation along two disjoint paths in the
mobility network and uses products formed from separate predictions. Where
the network cannot supply the paths or products it needs, it substitutes a
conservative variance contribution. Its published application reports inference
for firm variance. The model-based procedure in the KSS Matlab package instead
smooths raw leave-out products over leverage $P_{ii}$ and the component's
own weight $B_{ii}$. It fits a separate variance model for each of three
components. \fevc{} uses one set of variance inputs for all four components and their
covariance. Agreement between \fevc{} and the KSS Matlab package is agreement between
two model-based procedures, and both can err if the characteristics they pool over
miss variation in a few influential workers' wages. Agreement with the split-sample
procedure is the stronger check, because that procedure uses no variance model.

\subsection{Firm variance on the six Table IV samples}\label{kss:sec:emp-se}
The six samples of KSS Table~IV pair wages in 1999 and 2001 for pooled,
younger and older Veneto workers. Each group supplies a leave-one-out and
a more tightly connected leave-two-out sample. They were rebuilt from the raw data
with KSS's own code, and all \fevc{} runs in this section use observation deletion,
as KSS do, with the 2001 indicator as the only control. On common rows, with KSS's
year-effect treatment and enough projections, \fevc{} reproduces the KSS
firm-variance point estimates within numerical uncertainty. With two
observations per worker, leaving out one wage leaves the worker effect
pinned down by the other. The two leave-out products of a mover then sum
to the first-difference product used by KSS. The default estimates can
differ slightly because \fevc{} jointly estimates the year effect and approximates
leverages by random projections.

\begin{table}[htbp]\centering
\caption{Firm variance on KSS's Table IV samples: standard errors}\label{tab:six}
\small\begin{tabular}{@{}llrrrr@{}}
\toprule
Sample & Connected set & KSS split & KSS Matlab & \fevc{} & \fevc{} interval\\
\midrule
Pooled & Leave-one-out & .0006$^a$ & .00048 & .00049 & Normal\\
       & Leave-two-out & .0006$^a$ & .00053 & .00052 & Normal\\
Younger & Leave-one-out & .00059 & .00049 & .00052 & Normal\\
        & Leave-two-out & .00058 & .00055 & .00059 & $q=1$\\
Older  & Leave-one-out & .00244 & .00196 & .00199 & $q=1$\\
       & Leave-two-out & .00137 & .00110 & .00094 & $q=1$\\
\bottomrule
\end{tabular}
\par\smallskip\begin{minipage}{0.92\linewidth}\footnotesize
Entries are standard errors of the normal form, the square root of the variance
estimate. The final column identifies the interval reported by the \fevc{} default. In
the three $q=1$ rows the default prints the interval and no standard error; the entry
is the normal-form standard error kept in \code{e(inference\_intervals)}. KSS split
uses the Econometrica replication code, and KSS Matlab uses the model-based procedure
in commit \code{8b957ffe}. $^a$Published KSS values (at their printed precision).
The pooled split-sample reruns were stopped after 23 hours before the first
call finished.
\end{minipage}
\end{table}

Four findings stand out in Table~\ref{tab:six}.
\begin{itemize}
\item \emph{The two model-based standard errors agree.} On five of the six samples
they are within about 7\% of each other.
\item \emph{The interval choice lines up with KSS's intervals.} \fevc{} reports the $q=1$ interval
on exactly the three samples where KSS's $q=1$ and normal intervals differ, and the
normal interval on the three where they coincide.
\item \emph{Where the standard errors differ most, the intervals agree.} On the older
leave-two-out sample the standard errors range from .00094 for \fevc{} to .00137 for
KSS's split-sample calculation. Yet the reported $q=1$ intervals are almost identical:
$[.0154,.0218]$ for \fevc{} and $[.0155,.0220]$ for KSS. Those intervals depend on the
joint uncertainty of the leading score and the remainder, which a normal standard
error alone cannot summarise.
\item \emph{The split-sample standard errors are the largest where they are
conservative by construction.} On the leave-one-out samples, 15--22\% of movers lack
an independent second path and receive KSS's conservative substitution.
\end{itemize}
For worker variance and the covariance, for which KSS publish no standard errors,
the only comparison is the KSS Matlab package. The two agree within 6\% on the pooled
samples and within 11\% on the older samples, with no consistent sign. On the younger
samples the \fevc{} standard errors are 12--25\% larger.

\emph{Time.} On these samples \fevc{} took 18--96 seconds for all four components and
their covariance. The KSS Matlab package took roughly ten minutes to ten hours for
three components. The split-sample Table~IV driver took about thirteen minutes on the
older leave-two-out sample for its two calls on firm variance, and the pooled runs
were stopped after 23 hours. These are observed times on comparable cluster
nodes with different speeds, worker counts, leverage calculations and outputs.
They indicate orders of magnitude. The shared variance fit and the trace
calculation in \eqref{eq:sim} explain part of the difference.

\subsection{Against a known truth}\label{kss:sec:synth}
On real data there is no benchmark that determines which standard error is
closer. A simulation on the two older-worker networks, where the methods differ most,
preserves their actual moves and draws new wages from a version of KSS's Monte Carlo
design. Stayers have no wage noise, and each mover has one observation per employer,
so the simulation is a direct test of the fitted mover model. The design makes the
variance depend partly on the mover's weight $B_{ii}$ in the firm-variance estimator.
That characteristic is an input of the KSS Matlab smoother and is not in the model of
\fevc{}, so the design favours the former. The parameters were fitted on each
older-worker sample; KSS fitted theirs on the pooled sample.

\emph{On the leave-one-out network \fevc{} is on target.} Its firm-variance standard
error and that of the KSS Matlab package are within 4\% of the simulated spread, while
the split-sample standard error is 44\% too large and its intervals are 30\% longer.
The \fevc{} standard errors for worker variance and the covariance are 1\% and 4\%
below the spread, and its reported intervals cover about 97\% for all three
components.

\begin{table}[htbp]\centering
\caption{Older leave-two-out network: standard errors and coverage when the truth is known}
\label{tab:known}
\small\setlength{\tabcolsep}{4pt}\begin{tabular}{@{}llrrr@{}}
\toprule
Component & Procedure & Mean SE bias & Normal coverage & Reported coverage\\
\midrule
Firm variance & KSS split-sample & $+10\%$ & 95.5\% & 97.5\%\\
              & KSS Matlab model-based & $-5\%$ & 88.0\% & 76.5\%\\
              & \fevc{} default & $-6\%$ & 92.0\% & 98.9\%\\
\addlinespace
Worker variance & KSS Matlab model-based & $+2\%$ & 93.5\% & 95.5\%\\
                & \fevc{} default & $-14\%$ & 89.0\% & 97.6\%\\
Worker--firm covariance & KSS Matlab model-based & $+1\%$ & 90.0\% & 96.0\%\\
                        & \fevc{} default & $-15\%$ & 88.1\% & 97.5\%\\
\bottomrule
\end{tabular}
\par\smallskip\begin{minipage}{0.92\linewidth}\footnotesize
All simulated mover noise falls on the second year's wage in this design.
Bias compares the mean normal-form standard error with the standard
deviation across 1,000 \fevc{} estimates, net of numerical projection noise.
Normal coverage is that of the normal interval, which the \fevc{} default does not
report here. Reported coverage uses the selected interval, counting missing or invalid
intervals as misses; the KSS procedures use $q=1$ here. KSS procedures used the first 200
draws, so their coverage rates are less precise than those from 1,000.
\end{minipage}
\end{table}

\emph{On the leave-two-out network the reported intervals cover.} One direction
dominates there. For firm variance \fevc{} reported the $q=1$ interval on every draw,
and that interval has the same median length as the split-sample interval,
0.0086, with coverage of 98.9\% against 97.5\% (Table~\ref{tab:known}). The KSS Matlab
procedure produced shorter intervals (median length 0.0076), but 19\% of its $q=1$
intervals failed the study's validity criterion and counted as misses
(Appendix~\ref{app:evidence} gives the scoring convention). For worker variance and
the covariance the \fevc{} intervals covered 97.6\% and 97.5\%, and 96.1\% and 96.6\%
when the simulated noise is split across the two years.

\emph{The normal-form standard errors fall short on this network.} For firm variance
the shortfall is 6\%. For worker variance and the covariance it is 14--15\%, and
16--17\% with split noise, and normal intervals would cover only 88--89\%. As in the
noisy-bridge designs of Section~\ref{sec:q1sim}, the fitted model is wrong on exactly
the rows that carry the component. This is the situation in which the default does not
report a normal interval.

\emph{The shortfall traces to about thirty movers.} A single direction accounts for
95\% of the sampling variance of estimated worker variance and 92\% for the
covariance. The ten movers with the largest weights carry half the weight of that
direction. Their fitted error variances are just 44--48\% of their true values. Most
of the direction's weight is in leverage deciles six through nine, so inspecting only
the highest-leverage observations would miss it. Fitting the same model directly to
the true variances gives nearly the same discrepancy along the direction: the
functional form, not sampling noise, explains the shortfall. The larger-of-two rule of
Section~\ref{sec:q1int} is what protects the reported interval, because KSS's
leave-out variance of the leading score does not depend on the model.

This example gives content to $b(\Delta)$ in \eqref{eq:b}. An error in the
model matters to the extent that it falls on the workers who carry a
particular component. Whether actual wage noise rises on these links is
unknown. The coefficient on the firm-variance weight in the older sample
has the opposite sign from the estimate on KSS's pooled sample, and the
regression used to calibrate noise explains only 0.9\% of the variation
in the raw products. The simulation is therefore best read as a stress test: it shows
what happens when wage noise is tied to the very weights that carry a component, and
that the reported interval holds up even then.
